%% kapa: exam=1 round=30 year=2019 period=1 form=A subject=real_estate subject_ko=부동산학원론 seq=16 number=96
%% source: Q-Net 감정평가사 기출문제 (cst00309) artlSeq=5205205 file=extracted/1008/5205205/2019년 제30회 감정평가사 1차시험 문제지/2019년 제30회 감정평가사 1차시험 문제지/2019년도 제30회 감정평가사 1차 1교시 A형.hwp
%% conversion: kordoc 4.24.51 -> tools/gen1.py (auto); review=proofread-sonnet; flags=none
\begin{question}{96}
주택자금대출을 위한 다음과 같은 대안에 관한 설명으로 옳은 것은? (단, 주어진 조건에 한함) \par
{\small\centering
\begin{tabular}{|c|l|}
\hline
공통 대출조건 & \makecell[l]{○ 대출금: 2억원 \\ ○ 이자율: 고정금리, 연 5 \% \\ ○ 대출기간: 15년} \\
\hline
대안1 & 원금만기일시상환조건(이자는 연말납입) \\
\hline
대안2 & 원리금균등분할상환조건(거치기간 없이 연말상환) \\
\hline
대안3 & 원금균등분할상환조건(거치기간 없이 연말상환) \\
\hline
대안4 & 부(-)의 상환인 체증분할상환조건(거치기간 없이 연말상환) \\
\hline
\end{tabular}\par}

\begin{choices}
  \item 모든 대안별 대출금액에 대한 상환방식은 다르지만, 첫째년도에 지불하는 이자금액은 모든 대안이 동일하다.
  \item 모든 대안의 대출기간 동안에 상환한 원금과 이자의 총합계액은 동일하다.
  \item 대안4는 다른 대안에 비해서 대출기간이 경과할수록 이자부담이 점증하는 구조이기에 원금부담은 줄어든다.
  \item 대안2는 대안3에 비해서 첫째년도의 원금상환액이 큰 방식이다.
  \item 대안3은 다른 대안에 비해서 첫째년도에 차입자의 원리금지급 부담이 큰 방식이다.
\end{choices}
\end{question}
\answer{5}
